% Options for packages loaded elsewhere
\PassOptionsToPackage{unicode}{hyperref}
\PassOptionsToPackage{hyphens}{url}
\documentclass[
]{article}
\usepackage{xcolor}
\usepackage[margin=1in]{geometry}
\usepackage{amsmath,amssymb}
\setcounter{secnumdepth}{-\maxdimen} % remove section numbering
\usepackage{iftex}
\ifPDFTeX
  \usepackage[T1]{fontenc}
  \usepackage[utf8]{inputenc}
  \usepackage{textcomp} % provide euro and other symbols
\else % if luatex or xetex
  \usepackage{unicode-math} % this also loads fontspec
  \defaultfontfeatures{Scale=MatchLowercase}
  \defaultfontfeatures[\rmfamily]{Ligatures=TeX,Scale=1}
\fi
\usepackage{lmodern}
\ifPDFTeX\else
  % xetex/luatex font selection
\fi
% Use upquote if available, for straight quotes in verbatim environments
\IfFileExists{upquote.sty}{\usepackage{upquote}}{}
\IfFileExists{microtype.sty}{% use microtype if available
  \usepackage[]{microtype}
  \UseMicrotypeSet[protrusion]{basicmath} % disable protrusion for tt fonts
}{}
\makeatletter
\@ifundefined{KOMAClassName}{% if non-KOMA class
  \IfFileExists{parskip.sty}{%
    \usepackage{parskip}
  }{% else
    \setlength{\parindent}{0pt}
    \setlength{\parskip}{6pt plus 2pt minus 1pt}}
}{% if KOMA class
  \KOMAoptions{parskip=half}}
\makeatother
\usepackage{color}
\usepackage{fancyvrb}
\newcommand{\VerbBar}{|}
\newcommand{\VERB}{\Verb[commandchars=\\\{\}]}
\DefineVerbatimEnvironment{Highlighting}{Verbatim}{commandchars=\\\{\}}
% Add ',fontsize=\small' for more characters per line
\usepackage{framed}
\definecolor{shadecolor}{RGB}{248,248,248}
\newenvironment{Shaded}{\begin{snugshade}}{\end{snugshade}}
\newcommand{\AlertTok}[1]{\textcolor[rgb]{0.94,0.16,0.16}{#1}}
\newcommand{\AnnotationTok}[1]{\textcolor[rgb]{0.56,0.35,0.01}{\textbf{\textit{#1}}}}
\newcommand{\AttributeTok}[1]{\textcolor[rgb]{0.13,0.29,0.53}{#1}}
\newcommand{\BaseNTok}[1]{\textcolor[rgb]{0.00,0.00,0.81}{#1}}
\newcommand{\BuiltInTok}[1]{#1}
\newcommand{\CharTok}[1]{\textcolor[rgb]{0.31,0.60,0.02}{#1}}
\newcommand{\CommentTok}[1]{\textcolor[rgb]{0.56,0.35,0.01}{\textit{#1}}}
\newcommand{\CommentVarTok}[1]{\textcolor[rgb]{0.56,0.35,0.01}{\textbf{\textit{#1}}}}
\newcommand{\ConstantTok}[1]{\textcolor[rgb]{0.56,0.35,0.01}{#1}}
\newcommand{\ControlFlowTok}[1]{\textcolor[rgb]{0.13,0.29,0.53}{\textbf{#1}}}
\newcommand{\DataTypeTok}[1]{\textcolor[rgb]{0.13,0.29,0.53}{#1}}
\newcommand{\DecValTok}[1]{\textcolor[rgb]{0.00,0.00,0.81}{#1}}
\newcommand{\DocumentationTok}[1]{\textcolor[rgb]{0.56,0.35,0.01}{\textbf{\textit{#1}}}}
\newcommand{\ErrorTok}[1]{\textcolor[rgb]{0.64,0.00,0.00}{\textbf{#1}}}
\newcommand{\ExtensionTok}[1]{#1}
\newcommand{\FloatTok}[1]{\textcolor[rgb]{0.00,0.00,0.81}{#1}}
\newcommand{\FunctionTok}[1]{\textcolor[rgb]{0.13,0.29,0.53}{\textbf{#1}}}
\newcommand{\ImportTok}[1]{#1}
\newcommand{\InformationTok}[1]{\textcolor[rgb]{0.56,0.35,0.01}{\textbf{\textit{#1}}}}
\newcommand{\KeywordTok}[1]{\textcolor[rgb]{0.13,0.29,0.53}{\textbf{#1}}}
\newcommand{\NormalTok}[1]{#1}
\newcommand{\OperatorTok}[1]{\textcolor[rgb]{0.81,0.36,0.00}{\textbf{#1}}}
\newcommand{\OtherTok}[1]{\textcolor[rgb]{0.56,0.35,0.01}{#1}}
\newcommand{\PreprocessorTok}[1]{\textcolor[rgb]{0.56,0.35,0.01}{\textit{#1}}}
\newcommand{\RegionMarkerTok}[1]{#1}
\newcommand{\SpecialCharTok}[1]{\textcolor[rgb]{0.81,0.36,0.00}{\textbf{#1}}}
\newcommand{\SpecialStringTok}[1]{\textcolor[rgb]{0.31,0.60,0.02}{#1}}
\newcommand{\StringTok}[1]{\textcolor[rgb]{0.31,0.60,0.02}{#1}}
\newcommand{\VariableTok}[1]{\textcolor[rgb]{0.00,0.00,0.00}{#1}}
\newcommand{\VerbatimStringTok}[1]{\textcolor[rgb]{0.31,0.60,0.02}{#1}}
\newcommand{\WarningTok}[1]{\textcolor[rgb]{0.56,0.35,0.01}{\textbf{\textit{#1}}}}
\usepackage{graphicx}
\makeatletter
\newsavebox\pandoc@box
\newcommand*\pandocbounded[1]{% scales image to fit in text height/width
  \sbox\pandoc@box{#1}%
  \Gscale@div\@tempa{\textheight}{\dimexpr\ht\pandoc@box+\dp\pandoc@box\relax}%
  \Gscale@div\@tempb{\linewidth}{\wd\pandoc@box}%
  \ifdim\@tempb\p@<\@tempa\p@\let\@tempa\@tempb\fi% select the smaller of both
  \ifdim\@tempa\p@<\p@\scalebox{\@tempa}{\usebox\pandoc@box}%
  \else\usebox{\pandoc@box}%
  \fi%
}
% Set default figure placement to htbp
\def\fps@figure{htbp}
\makeatother
\setlength{\emergencystretch}{3em} % prevent overfull lines
\providecommand{\tightlist}{%
  \setlength{\itemsep}{0pt}\setlength{\parskip}{0pt}}
\usepackage{bookmark}
\IfFileExists{xurl.sty}{\usepackage{xurl}}{} % add URL line breaks if available
\urlstyle{same}
\hypersetup{
  pdftitle={HW\_04},
  pdfauthor={Henry Traynor},
  hidelinks,
  pdfcreator={LaTeX via pandoc}}

\title{HW\_04}
\author{Henry Traynor}
\date{Due: 5:00pm, 15 September}

\begin{document}
\maketitle

Create a self-contained R script or Rmarkdown/Quarto file to do the
following:

Using the simulated data, compare the prior and posterior predictive
distributions of \(\lambda\) for each of the 3 forest types.

Hint: because forest type is categorical (not continuous), you can
visualize the prior and posterior predictive distributions as histograms
for each forest type.

\begin{Shaded}
\begin{Highlighting}[]
\FunctionTok{set.seed}\NormalTok{(}\DecValTok{314159}\NormalTok{) }\CommentTok{\# for reproducibility}

\NormalTok{nSites }\OtherTok{\textless{}{-}} \DecValTok{100}
\NormalTok{nVisits }\OtherTok{\textless{}{-}} \DecValTok{4}

\NormalTok{forest }\OtherTok{\textless{}{-}} \FunctionTok{factor}\NormalTok{(}\FunctionTok{sample}\NormalTok{(}\FunctionTok{c}\NormalTok{(}\StringTok{"Hardwood"}\NormalTok{, }\StringTok{"Mixed"}\NormalTok{, }\StringTok{"Pine"}\NormalTok{), nSites, }\AttributeTok{replace=}\ConstantTok{TRUE}\NormalTok{))}
\NormalTok{forestMixed }\OtherTok{\textless{}{-}} \FunctionTok{ifelse}\NormalTok{(forest}\SpecialCharTok{==}\StringTok{"Mixed"}\NormalTok{, }\DecValTok{1}\NormalTok{, }\DecValTok{0}\NormalTok{)        }\DocumentationTok{\#\# Dummy}
\NormalTok{forestPine }\OtherTok{\textless{}{-}} \FunctionTok{ifelse}\NormalTok{(forest}\SpecialCharTok{==}\StringTok{"Pine"}\NormalTok{, }\DecValTok{1}\NormalTok{, }\DecValTok{0}\NormalTok{)          }\DocumentationTok{\#\# Dummy}
\NormalTok{temp }\OtherTok{\textless{}{-}} \FunctionTok{matrix}\NormalTok{(}\FunctionTok{rnorm}\NormalTok{(nSites}\SpecialCharTok{*}\NormalTok{nVisits), }\AttributeTok{nrow=}\NormalTok{nSites, }\AttributeTok{ncol=}\NormalTok{nVisits)}

\CommentTok{\# true coeffs}
\NormalTok{beta0 }\OtherTok{\textless{}{-}} \DecValTok{0}
\NormalTok{beta1 }\OtherTok{\textless{}{-}} \SpecialCharTok{{-}}\DecValTok{1}
\NormalTok{beta2 }\OtherTok{\textless{}{-}} \DecValTok{1}
\NormalTok{lambda }\OtherTok{\textless{}{-}} \FunctionTok{exp}\NormalTok{(beta0 }\SpecialCharTok{+}\NormalTok{ beta1}\SpecialCharTok{*}\NormalTok{forestMixed }\SpecialCharTok{+}\NormalTok{ beta2}\SpecialCharTok{*}\NormalTok{forestPine)}
\NormalTok{alpha0 }\OtherTok{\textless{}{-}} \SpecialCharTok{{-}}\DecValTok{2}
\NormalTok{alpha1 }\OtherTok{\textless{}{-}} \DecValTok{1}
\NormalTok{p }\OtherTok{\textless{}{-}} \FunctionTok{plogis}\NormalTok{(alpha0 }\SpecialCharTok{+}\NormalTok{ alpha1}\SpecialCharTok{*}\NormalTok{temp)}

\CommentTok{\# true lambda plots}
\NormalTok{lambda.true }\OtherTok{\textless{}{-}} \FunctionTok{c}\NormalTok{(}
  \AttributeTok{Hardwood =} \FunctionTok{exp}\NormalTok{(beta0),}
  \AttributeTok{Mixed =} \FunctionTok{exp}\NormalTok{(beta0 }\SpecialCharTok{+}\NormalTok{ beta1),}
  \AttributeTok{Pine =} \FunctionTok{exp}\NormalTok{(beta0 }\SpecialCharTok{+}\NormalTok{ beta2)}
\NormalTok{)}
\NormalTok{lambda.true}
\end{Highlighting}
\end{Shaded}

\begin{verbatim}
##  Hardwood     Mixed      Pine 
## 1.0000000 0.3678794 2.7182818
\end{verbatim}

\begin{Shaded}
\begin{Highlighting}[]
\CommentTok{\# simulate occupancy and detection data}
\NormalTok{N }\OtherTok{\textless{}{-}} \FunctionTok{rpois}\NormalTok{(nSites, }\AttributeTok{lambda=}\NormalTok{lambda)         }\DocumentationTok{\#\# local abundance }
\NormalTok{y }\OtherTok{\textless{}{-}} \FunctionTok{matrix}\NormalTok{(}\ConstantTok{NA}\NormalTok{, }\AttributeTok{nrow=}\NormalTok{nSites, }\AttributeTok{ncol=}\NormalTok{nVisits)}
\ControlFlowTok{for}\NormalTok{(i }\ControlFlowTok{in} \DecValTok{1}\SpecialCharTok{:}\NormalTok{nSites) \{}
\NormalTok{    y[i,] }\OtherTok{\textless{}{-}} \FunctionTok{rbinom}\NormalTok{(nVisits, }\AttributeTok{size=}\NormalTok{N[i], }\AttributeTok{prob=}\NormalTok{p[i,])}
\NormalTok{\}}

\CommentTok{\# summary statistics}
\NormalTok{maxCounts }\OtherTok{\textless{}{-}} \FunctionTok{apply}\NormalTok{(y, }\DecValTok{1}\NormalTok{, max) }
\FunctionTok{table}\NormalTok{(maxCounts)            }
\end{Highlighting}
\end{Shaded}

\begin{verbatim}
## maxCounts
##  0  1  2  3  4 
## 61 34  2  2  1
\end{verbatim}

\begin{Shaded}
\begin{Highlighting}[]
\NormalTok{naiveOccupancy }\OtherTok{\textless{}{-}} \FunctionTok{sum}\NormalTok{(maxCounts}\SpecialCharTok{\textgreater{}}\DecValTok{0}\NormalTok{)}\SpecialCharTok{/}\NormalTok{nSites}
\NormalTok{naiveOccupancy }
\end{Highlighting}
\end{Shaded}

\begin{verbatim}
## [1] 0.39
\end{verbatim}

\begin{Shaded}
\begin{Highlighting}[]
\NormalTok{naiveAbund }\OtherTok{\textless{}{-}} \FunctionTok{sum}\NormalTok{(maxCounts)}
\NormalTok{naiveAbund}
\end{Highlighting}
\end{Shaded}

\begin{verbatim}
## [1] 48
\end{verbatim}

\begin{Shaded}
\begin{Highlighting}[]
\CommentTok{\# prep data in named list}
\NormalTok{jags.data }\OtherTok{\textless{}{-}} \FunctionTok{list}\NormalTok{(}\AttributeTok{y=}\NormalTok{y,}
                  \AttributeTok{forestMixed=}\NormalTok{forestMixed,}
                  \AttributeTok{forestPine=}\NormalTok{forestPine,}
                  \AttributeTok{nSites=}\NormalTok{nSites, }
                  \AttributeTok{temp=}\NormalTok{temp,}
                  \AttributeTok{nOccasions=}\NormalTok{nVisits)}
\CommentTok{\# initial values for MCMC}
\NormalTok{jags.inits }\OtherTok{\textless{}{-}} \ControlFlowTok{function}\NormalTok{() \{}
    \FunctionTok{list}\NormalTok{(}\AttributeTok{beta0=}\FunctionTok{rnorm}\NormalTok{(}\DecValTok{1}\NormalTok{), }\AttributeTok{alpha0=}\FunctionTok{rnorm}\NormalTok{(}\DecValTok{1}\NormalTok{), }\AttributeTok{N=}\NormalTok{maxCounts)}
\NormalTok{\}}
\CommentTok{\# parameters to save}
\NormalTok{jags.pars }\OtherTok{\textless{}{-}} \FunctionTok{c}\NormalTok{(}\StringTok{"beta0"}\NormalTok{, }\StringTok{"beta1"}\NormalTok{, }\StringTok{"beta2"}\NormalTok{,}
               \StringTok{"alpha0"}\NormalTok{, }\StringTok{"alpha1"}\NormalTok{, }\StringTok{"totalAbundance"}\NormalTok{)}

\CommentTok{\# run model}
\FunctionTok{library}\NormalTok{(jagsUI)}
\NormalTok{jags.post.samples }\OtherTok{\textless{}{-}} \FunctionTok{jags.basic}\NormalTok{(}\AttributeTok{data=}\NormalTok{jags.data, }
                                \AttributeTok{inits=}\NormalTok{jags.inits,}
                                \AttributeTok{parameters.to.save=}\FunctionTok{c}\NormalTok{(jags.pars,}\StringTok{"N"}\NormalTok{), }
                                \AttributeTok{model.file=}\StringTok{"Nmix{-}model{-}covs.jag"}\NormalTok{,}
                                \AttributeTok{n.chains=}\DecValTok{3}\NormalTok{,}
                                \AttributeTok{n.adapt=}\DecValTok{100}\NormalTok{, }
                                \AttributeTok{n.burnin=}\DecValTok{0}\NormalTok{,}
                                \AttributeTok{n.iter=}\DecValTok{6000}\NormalTok{,}
                                \AttributeTok{parallel=}\ConstantTok{TRUE}\NormalTok{)}
\end{Highlighting}
\end{Shaded}

\begin{verbatim}
## 
## Processing function input....... 
## 
## Done. 
##  
## Beginning parallel processing using 3 cores. Console output will be suppressed.
## 
## Parallel processing completed.
## 
## MCMC took 0.154 minutes.
\end{verbatim}

\begin{Shaded}
\begin{Highlighting}[]
\FunctionTok{summary}\NormalTok{(jags.post.samples[,jags.pars])}
\end{Highlighting}
\end{Shaded}

\begin{verbatim}
## 
## Iterations = 101:6100
## Thinning interval = 1 
## Number of chains = 3 
## Sample size per chain = 6000 
## 
## 1. Empirical mean and standard deviation for each variable,
##    plus standard error of the mean:
## 
##                   Mean      SD Naive SE Time-series SE
## beta0          -0.2160  0.2736 0.002039        0.01795
## beta1          -0.7590  0.4029 0.003003        0.02311
## beta2           0.8277  0.3097 0.002308        0.01427
## alpha0         -1.6143  0.3127 0.002330        0.04394
## alpha1          1.0628  0.1735 0.001293        0.01532
## totalAbundance 94.2259 20.2886 0.151222        2.70765
## 
## 2. Quantiles for each variable:
## 
##                   2.5%     25%     50%       75%      97.5%
## beta0          -0.7423 -0.4023 -0.2243  -0.03539   0.344070
## beta1          -1.5837 -1.0244 -0.7474  -0.48342   0.005029
## beta2           0.2328  0.6187  0.8249   1.03222   1.443433
## alpha0         -2.2454 -1.8118 -1.5926  -1.40132  -1.018052
## alpha1          0.7235  0.9499  1.0558   1.17456   1.408526
## totalAbundance 66.0000 80.0000 90.0000 104.00000 145.000000
\end{verbatim}

\begin{Shaded}
\begin{Highlighting}[]
\FunctionTok{plot}\NormalTok{(jags.post.samples[,}\FunctionTok{paste0}\NormalTok{(}\StringTok{"N["}\NormalTok{, }\DecValTok{1}\SpecialCharTok{:}\DecValTok{2}\NormalTok{, }\StringTok{"]"}\NormalTok{)])}
\end{Highlighting}
\end{Shaded}

\pandocbounded{\includegraphics[keepaspectratio]{HW_04_files/HW_04_files/figure-latex/Summary & Plots-1.pdf}}

\begin{Shaded}
\begin{Highlighting}[]
\FunctionTok{plot}\NormalTok{(jags.post.samples[,}\FunctionTok{paste0}\NormalTok{(}\StringTok{"N["}\NormalTok{, }\DecValTok{3}\SpecialCharTok{:}\DecValTok{4}\NormalTok{, }\StringTok{"]"}\NormalTok{)])}
\end{Highlighting}
\end{Shaded}

\pandocbounded{\includegraphics[keepaspectratio]{HW_04_files/HW_04_files/figure-latex/Summary & Plots-2.pdf}}

\begin{Shaded}
\begin{Highlighting}[]
\FunctionTok{plot}\NormalTok{(jags.post.samples[,jags.pars[}\DecValTok{1}\SpecialCharTok{:}\DecValTok{3}\NormalTok{]])}
\end{Highlighting}
\end{Shaded}

\pandocbounded{\includegraphics[keepaspectratio]{HW_04_files/HW_04_files/figure-latex/Summary & Plots-3.pdf}}

\begin{Shaded}
\begin{Highlighting}[]
\FunctionTok{plot}\NormalTok{(jags.post.samples[,jags.pars[}\DecValTok{4}\SpecialCharTok{:}\DecValTok{5}\NormalTok{]])}
\end{Highlighting}
\end{Shaded}

\pandocbounded{\includegraphics[keepaspectratio]{HW_04_files/HW_04_files/figure-latex/Summary & Plots-4.pdf}}

\begin{Shaded}
\begin{Highlighting}[]
\CommentTok{\# extract beta values}
\NormalTok{lambda.coef.post }\OtherTok{\textless{}{-}} \FunctionTok{as.matrix}\NormalTok{(jags.post.samples[,}\FunctionTok{c}\NormalTok{(}\StringTok{"beta0"}\NormalTok{,}\StringTok{"beta1"}\NormalTok{, }\StringTok{"beta2"}\NormalTok{)])}
\FunctionTok{head}\NormalTok{(lambda.coef.post, }\AttributeTok{n=}\DecValTok{4}\NormalTok{)}
\end{Highlighting}
\end{Shaded}

\begin{verbatim}
##           beta0      beta1     beta2
## [1,] -0.4837185 -0.7557389 0.3601422
## [2,] -0.4710857 -0.6743193 0.3341998
## [3,] -0.3907293 -0.7037765 0.3400097
## [4,] -0.3746024 -0.7438384 0.4104394
\end{verbatim}

\begin{Shaded}
\begin{Highlighting}[]
\NormalTok{n.iter }\OtherTok{\textless{}{-}} \FunctionTok{nrow}\NormalTok{(lambda.coef.post)}
\CommentTok{\# vector for predictive val}
\NormalTok{forest.pred }\OtherTok{\textless{}{-}} \FunctionTok{c}\NormalTok{(}\FunctionTok{rep}\NormalTok{(}\StringTok{"Hardwood"}\NormalTok{, }\DecValTok{50}\NormalTok{),}
                 \FunctionTok{rep}\NormalTok{(}\StringTok{"Mixed"}\NormalTok{, }\DecValTok{50}\NormalTok{),}
                 \FunctionTok{rep}\NormalTok{(}\StringTok{"Pine"}\NormalTok{,}\DecValTok{50}\NormalTok{)}
\NormalTok{                 )}

\NormalTok{forest.Mixed.pred }\OtherTok{\textless{}{-}} \FunctionTok{ifelse}\NormalTok{(forest.pred}\SpecialCharTok{==}\StringTok{"Mixed"}\NormalTok{, }\DecValTok{1}\NormalTok{, }\DecValTok{0}\NormalTok{)}
\NormalTok{forest.Pine.pred }\OtherTok{\textless{}{-}} \FunctionTok{ifelse}\NormalTok{(forest.pred}\SpecialCharTok{==}\StringTok{"Pine"}\NormalTok{, }\DecValTok{1}\NormalTok{, }\DecValTok{0}\NormalTok{)}

\NormalTok{lambda.post.pred }\OtherTok{\textless{}{-}} \FunctionTok{matrix}\NormalTok{(}\ConstantTok{NA}\NormalTok{, }\AttributeTok{nrow=}\NormalTok{n.iter, }\AttributeTok{ncol=}\FunctionTok{length}\NormalTok{(forest.pred))}

\CommentTok{\# we now use exp() rather than plogis() since lambda is log{-}linear rather than logit{-}linear}
\ControlFlowTok{for}\NormalTok{(i }\ControlFlowTok{in} \DecValTok{1}\SpecialCharTok{:}\NormalTok{n.iter) \{}
\NormalTok{    lambda.post.pred[i,] }\OtherTok{\textless{}{-}} \FunctionTok{exp}\NormalTok{(}
\NormalTok{      lambda.coef.post[i,}\StringTok{"beta0"}\NormalTok{] }\SpecialCharTok{+} 
\NormalTok{      lambda.coef.post[i,}\StringTok{"beta1"}\NormalTok{] }\SpecialCharTok{*}\NormalTok{ forest.Mixed.pred }\SpecialCharTok{+} 
\NormalTok{      lambda.coef.post[i,}\StringTok{"beta2"}\NormalTok{] }\SpecialCharTok{*}\NormalTok{ forest.Pine.pred}
\NormalTok{      )}
\NormalTok{\}}

\CommentTok{\# plot showing forest type on categorical horiz. axis and lambda on continuous vert. vertical axis}
\CommentTok{\# pull out values for each forest type}
\NormalTok{hardwood }\OtherTok{\textless{}{-}} \FunctionTok{as.vector}\NormalTok{(lambda.post.pred[, }\DecValTok{1}\SpecialCharTok{:}\DecValTok{50}\NormalTok{])}
\NormalTok{mixed    }\OtherTok{\textless{}{-}} \FunctionTok{as.vector}\NormalTok{(lambda.post.pred[, }\DecValTok{51}\SpecialCharTok{:}\DecValTok{100}\NormalTok{])}
\NormalTok{pine     }\OtherTok{\textless{}{-}} \FunctionTok{as.vector}\NormalTok{(lambda.post.pred[, }\DecValTok{101}\SpecialCharTok{:}\DecValTok{150}\NormalTok{])}
\CommentTok{\#combine vectors in long form for easy graphing}
\NormalTok{lambda.long }\OtherTok{\textless{}{-}} \FunctionTok{data.frame}\NormalTok{(}
  \AttributeTok{lambda =} \FunctionTok{c}\NormalTok{(hardwood, mixed, pine),}
  \AttributeTok{forest =} \FunctionTok{rep}\NormalTok{(}
    \FunctionTok{c}\NormalTok{(}\StringTok{"Hardwood"}\NormalTok{, }\StringTok{"Mixed"}\NormalTok{, }\StringTok{"Pine"}\NormalTok{),}
    \AttributeTok{times =} \FunctionTok{c}\NormalTok{(}\FunctionTok{length}\NormalTok{(hardwood), }\FunctionTok{length}\NormalTok{(mixed), }\FunctionTok{length}\NormalTok{(pine))}
\NormalTok{  )}
\NormalTok{)}
\FunctionTok{library}\NormalTok{(ggplot2)}
\FunctionTok{ggplot}\NormalTok{(lambda.long, }\FunctionTok{aes}\NormalTok{(}\AttributeTok{x =}\NormalTok{ lambda)) }\SpecialCharTok{+}
  \FunctionTok{geom\_histogram}\NormalTok{() }\SpecialCharTok{+}
  \FunctionTok{facet\_wrap}\NormalTok{(}\SpecialCharTok{\textasciitilde{}}\NormalTok{forest, }\AttributeTok{ncol =} \DecValTok{3}\NormalTok{) }\SpecialCharTok{+}
  \FunctionTok{labs}\NormalTok{(}
    \AttributeTok{x =} \StringTok{"lambda\_hat"}\NormalTok{,}
    \AttributeTok{y =} \StringTok{"Frequency"}
\NormalTok{  ) }\SpecialCharTok{+} 
  \FunctionTok{geom\_vline}\NormalTok{(}
    \AttributeTok{data =} \FunctionTok{data.frame}\NormalTok{(}
      \AttributeTok{forest =} \FunctionTok{names}\NormalTok{(lambda.true),}
      \AttributeTok{lambda =}\NormalTok{ lambda.true}
\NormalTok{    ),}
    \FunctionTok{aes}\NormalTok{(}\AttributeTok{xintercept =}\NormalTok{ lambda, }\AttributeTok{color =} \StringTok{"True value"}\NormalTok{),}
    \AttributeTok{linetype =} \StringTok{"solid"}
\NormalTok{)}
\end{Highlighting}
\end{Shaded}

\begin{verbatim}
## `stat_bin()` using `bins = 30`. Pick better value `binwidth`.
\end{verbatim}

\pandocbounded{\includegraphics[keepaspectratio]{HW_04_files/HW_04_files/figure-latex/Bayesian Prediction-1.pdf}}

\begin{Shaded}
\begin{Highlighting}[]
\NormalTok{jags.data2 }\OtherTok{\textless{}{-}}\NormalTok{ jags.data }\CommentTok{\# redefine data}
\NormalTok{jags.data2}\SpecialCharTok{$}\NormalTok{y[] }\OtherTok{\textless{}{-}} \ConstantTok{NA}     \CommentTok{\# Replace data with missing values}

\CommentTok{\# resample}
\NormalTok{jags.prior.samples }\OtherTok{\textless{}{-}} \FunctionTok{jags.basic}\NormalTok{(}\AttributeTok{data=}\NormalTok{jags.data2,}
                                 \AttributeTok{inits=}\NormalTok{jags.inits,}
                                 \AttributeTok{parameters.to.save=}\NormalTok{jags.pars,}
                                 \AttributeTok{model.file=}\StringTok{"Nmix{-}model{-}covs.jag"}\NormalTok{,}
                                 \AttributeTok{n.chains=}\DecValTok{3}\NormalTok{,}
                                 \AttributeTok{n.adapt=}\DecValTok{100}\NormalTok{,}
                                 \AttributeTok{n.burnin=}\DecValTok{0}\NormalTok{,}
                                 \AttributeTok{n.iter=}\DecValTok{6000}\NormalTok{,}
                                 \AttributeTok{parallel=}\ConstantTok{TRUE}\NormalTok{)}
\end{Highlighting}
\end{Shaded}

\begin{verbatim}
## 
## Processing function input....... 
## 
## Done. 
##  
## Beginning parallel processing using 3 cores. Console output will be suppressed.
## 
## Parallel processing completed.
## 
## MCMC took 0.01 minutes.
\end{verbatim}

\begin{Shaded}
\begin{Highlighting}[]
\FunctionTok{summary}\NormalTok{(jags.prior.samples)}
\end{Highlighting}
\end{Shaded}

\begin{verbatim}
## 
## Iterations = 1:6000
## Thinning interval = 1 
## Number of chains = 3 
## Sample size per chain = 6000 
## 
## 1. Empirical mean and standard deviation for each variable,
##    plus standard error of the mean:
## 
##                      Mean       SD Naive SE Time-series SE
## alpha0          -0.012572    1.422  0.01060        0.01060
## alpha1          -0.003566    1.412  0.01053        0.01029
## beta0           -0.006604    1.412  0.01052        0.01036
## beta1           -0.015821    1.435  0.01069        0.01069
## beta2            0.008797    1.393  0.01039        0.01032
## totalAbundance 547.307778 2365.565 17.63188       17.63236
## 
## 2. Quantiles for each variable:
## 
##                  2.5%     25%        50%      75%    97.5%
## alpha0         -2.779 -0.9815  -0.012361   0.9576    2.766
## alpha1         -2.785 -0.9442  -0.012058   0.9313    2.789
## beta0          -2.780 -0.9548  -0.013094   0.9527    2.767
## beta1          -2.815 -0.9968  -0.009533   0.9507    2.780
## beta2          -2.753 -0.9247   0.002908   0.9524    2.733
## totalAbundance  6.000 47.0000 137.000000 410.0000 3386.275
\end{verbatim}

\begin{Shaded}
\begin{Highlighting}[]
\NormalTok{lambda.coef.prior }\OtherTok{\textless{}{-}} \FunctionTok{as.matrix}\NormalTok{(jags.prior.samples[,}\FunctionTok{c}\NormalTok{(}\StringTok{"beta0"}\NormalTok{,}\StringTok{"beta1"}\NormalTok{,}\StringTok{"beta2"}\NormalTok{)])}
\NormalTok{lambda.prior.pred }\OtherTok{\textless{}{-}} \FunctionTok{matrix}\NormalTok{(}\ConstantTok{NA}\NormalTok{, }\AttributeTok{nrow=}\NormalTok{n.iter, }\AttributeTok{ncol=}\FunctionTok{length}\NormalTok{(forest.pred))}
\ControlFlowTok{for}\NormalTok{(i }\ControlFlowTok{in} \DecValTok{1}\SpecialCharTok{:}\NormalTok{n.iter) \{}
\NormalTok{    lambda.prior.pred[i,] }\OtherTok{\textless{}{-}} \FunctionTok{exp}\NormalTok{(lambda.coef.prior[i,}\StringTok{"beta0"}\NormalTok{] }\SpecialCharTok{+}\NormalTok{ lambda.coef.prior[i,}\StringTok{"beta1"}\NormalTok{]}\SpecialCharTok{*}\NormalTok{forest.Mixed.pred }\SpecialCharTok{+}\NormalTok{ lambda.coef.prior[i, }\StringTok{"beta2"}\NormalTok{]}\SpecialCharTok{*}\NormalTok{forest.Pine.pred)}
\NormalTok{\}}

\CommentTok{\# plot showing forest type on categorical horiz. axis and lambda on continuous vert. vertical axis}
\NormalTok{hardwood.prior }\OtherTok{\textless{}{-}} \FunctionTok{as.vector}\NormalTok{(lambda.prior.pred[, }\DecValTok{1}\SpecialCharTok{:}\DecValTok{50}\NormalTok{])}
\NormalTok{mixed.prior    }\OtherTok{\textless{}{-}} \FunctionTok{as.vector}\NormalTok{(lambda.prior.pred[, }\DecValTok{51}\SpecialCharTok{:}\DecValTok{100}\NormalTok{])}
\NormalTok{pine.prior     }\OtherTok{\textless{}{-}} \FunctionTok{as.vector}\NormalTok{(lambda.prior.pred[, }\DecValTok{101}\SpecialCharTok{:}\DecValTok{150}\NormalTok{])}
\NormalTok{lambda.long.prior }\OtherTok{\textless{}{-}} \FunctionTok{data.frame}\NormalTok{(}
  \AttributeTok{lambda =} \FunctionTok{c}\NormalTok{(hardwood.prior, mixed.prior, pine.prior),}
  \AttributeTok{forest =} \FunctionTok{rep}\NormalTok{(}
    \FunctionTok{c}\NormalTok{(}\StringTok{"Hardwood"}\NormalTok{, }\StringTok{"Mixed"}\NormalTok{, }\StringTok{"Pine"}\NormalTok{),}
    \AttributeTok{times =} \FunctionTok{c}\NormalTok{(}
      \FunctionTok{length}\NormalTok{(hardwood.prior),}
      \FunctionTok{length}\NormalTok{(mixed.prior),}
      \FunctionTok{length}\NormalTok{(pine.prior)}
\NormalTok{    )}
\NormalTok{  )}
\NormalTok{)}

\NormalTok{log.lambda.long.prior }\OtherTok{\textless{}{-}}\NormalTok{ lambda.long.prior}
\NormalTok{log.lambda.long.prior}\SpecialCharTok{$}\NormalTok{lambda }\OtherTok{\textless{}{-}} \FunctionTok{log}\NormalTok{(log.lambda.long.prior}\SpecialCharTok{$}\NormalTok{lambda)}
\FunctionTok{ggplot}\NormalTok{(log.lambda.long.prior, }\FunctionTok{aes}\NormalTok{(}\AttributeTok{x =}\NormalTok{ lambda)) }\SpecialCharTok{+}
  \FunctionTok{geom\_histogram}\NormalTok{() }\SpecialCharTok{+}
  \FunctionTok{facet\_wrap}\NormalTok{(}\SpecialCharTok{\textasciitilde{}}\NormalTok{forest, }\AttributeTok{ncol =} \DecValTok{3}\NormalTok{) }\SpecialCharTok{+}
  \FunctionTok{labs}\NormalTok{(}
    \AttributeTok{x =} \StringTok{"log(lambda\_hat)"}\NormalTok{,}
    \AttributeTok{y =} \StringTok{"Frequency"}
\NormalTok{  ) }\SpecialCharTok{+} 
  \FunctionTok{geom\_vline}\NormalTok{(}
    \AttributeTok{data =} \FunctionTok{data.frame}\NormalTok{(}
      \AttributeTok{forest =} \FunctionTok{names}\NormalTok{(lambda.true),}
      \AttributeTok{lambda =}\NormalTok{ lambda.true}
\NormalTok{    ),}
    \FunctionTok{aes}\NormalTok{(}\AttributeTok{xintercept =}\NormalTok{ lambda, }\AttributeTok{color =} \StringTok{"True value"}\NormalTok{),}
    \AttributeTok{linetype =} \StringTok{"solid"}
\NormalTok{  )}
\end{Highlighting}
\end{Shaded}

\begin{verbatim}
## `stat_bin()` using `bins = 30`. Pick better value `binwidth`.
\end{verbatim}

\pandocbounded{\includegraphics[keepaspectratio]{HW_04_files/HW_04_files/figure-latex/Prior Prediction-1.pdf}}

\begin{Shaded}
\begin{Highlighting}[]
\NormalTok{lambda.long}\SpecialCharTok{$}\NormalTok{distribution }\OtherTok{\textless{}{-}} \StringTok{"Posterior"}
\NormalTok{log.lambda.long.prior}\SpecialCharTok{$}\NormalTok{distribution }\OtherTok{\textless{}{-}} \StringTok{"Prior"}
\NormalTok{lambda.combined }\OtherTok{\textless{}{-}} \FunctionTok{rbind}\NormalTok{(lambda.long, log.lambda.long.prior)}

\FunctionTok{ggplot}\NormalTok{(lambda.combined,}
       \FunctionTok{aes}\NormalTok{(}\AttributeTok{x =}\NormalTok{ lambda, }\AttributeTok{fill =}\NormalTok{ distribution)) }\SpecialCharTok{+}
  \FunctionTok{geom\_histogram}\NormalTok{(}
    \AttributeTok{position =} \StringTok{"identity"}\NormalTok{,}
    \AttributeTok{alpha =} \FloatTok{0.5}\NormalTok{,}
\NormalTok{  ) }\SpecialCharTok{+}
  \FunctionTok{facet\_wrap}\NormalTok{(}\SpecialCharTok{\textasciitilde{}}\NormalTok{ forest, }\AttributeTok{ncol =} \DecValTok{3}\NormalTok{) }\SpecialCharTok{+}
  \FunctionTok{labs}\NormalTok{(}
    \AttributeTok{x =} \StringTok{"Posterior: lambda\_hat }\SpecialCharTok{\textbackslash{}n}\StringTok{ Prior: log(lambda\_hat)"}\NormalTok{,}
    \AttributeTok{y =} \StringTok{"Frequency"}\NormalTok{,}
    \AttributeTok{fill =} \StringTok{"Distribution"}
\NormalTok{  ) }\SpecialCharTok{+}
  \FunctionTok{geom\_vline}\NormalTok{(}
      \AttributeTok{data =} \FunctionTok{data.frame}\NormalTok{(}
        \AttributeTok{forest =} \FunctionTok{names}\NormalTok{(lambda.true),}
        \AttributeTok{lambda =}\NormalTok{ lambda.true}
\NormalTok{      ),}
      \FunctionTok{aes}\NormalTok{(}\AttributeTok{xintercept =}\NormalTok{ lambda, }\AttributeTok{color =} \StringTok{"True value, original scale"}\NormalTok{),}
      \AttributeTok{linetype =} \StringTok{"solid"}
\NormalTok{  )}
\end{Highlighting}
\end{Shaded}

\begin{verbatim}
## `stat_bin()` using `bins = 30`. Pick better value `binwidth`.
\end{verbatim}

\pandocbounded{\includegraphics[keepaspectratio]{HW_04_files/HW_04_files/figure-latex/Comparison Plot-1.pdf}}
Note:

I have showed the posterior \(\lambda\) distribution on it's original
scale. Due to the relatively much wider prior distribution, I have
displayed the prior \(\ambda\) distribution on the log-scale, as shown
in the horizontal axis label. The vertical line showing the true value
of \(\lambda\) is on the original scale and thus is only applicable when
visually expecting the accuracy of the posterior distribution, not the
prior distribution.

\end{document}
